\section*{1.1}
The goal is to analyse some high-dimensional system using low-dimensional models, using the hypothesis that some dimensions of the system are correlated. This chapter aims to explain some existing methods to find some intrisic low-dimensional structures. An example is the reconstruction of images, well-researched using SVD, where the columns of pixels are highly correlated and only a few allow to reconstruct the image quite well. 

This section aims to reduce high-dimensional dynamical systems to lower-dimensional ones, by applying some linear applications on the input and output of the system, allowing to use less observations. 

\section*{3.1.2}
As numerical methods can only consider discrete-time systems and not continuous ones, we need a way to discretize the systems correctly, in order to capture as much information as the continuous formulation. The low-dimensional hypothesis allows to a relatively small number of samples because of the low rank of the modifies output $YU^T$. 

This comes from the Fourier transform of the signal $x(t)$:
\begin{equation}
    x(t) = \frac{1}{2}\int_{-\infty}^\infty \hat x(\omega)\exp{(i\omega t)}d\omega = \frac{1}{2\pi} \int_{-\Omega}^\Omega \hat x(\omega) \exp{(i\omega t)}d\omega 
\end{equation}
If $\hat x(\omega)$ is viewed as a periodic function in the spectral domain with a period $2\Omega$, it is then fully determined by its Fourier coefficients: 
\begin{equation}
    x\left(\frac{n\pi}{\Omega}\right) = \frac{1}{2\pi}\int_{-\Omega}^{\Omega} \hat x(\omega) \exp{\left(i\omega \frac{n\pi}{\Omega}\right)}d\omega
\end{equation}
and $f=\frac{\Omega}{\pi}$ is the frequency of the sampling. 
\begin{itemize}
    \item [$\bullet$] Note: to perfectly recover a band-limited signal $x(t)$, we need to sample it at a rate that is twice its maximal frequency $f_{\max} = \Omega/2\pi$.
\end{itemize}

\section*{3.1.3}
One may view the observed data as samples of a set of random variables $\boldsymbol{x}_o\in \R^{n_o}$ generated through certain conditional probability distribution given another set of hidden (said latent) variables $\boldsymbol{x}_h\in \R^{n_h}$. 

In the case of jointly Gaussian variables, we can partition the covariance matrix as follows. 

meaning that is can be represented easily, compared to a full matrix of size $n_o\times n_o$. 

demander à Charles et pas PAA.

\section*{3.2}
